\section*{Sammendrag}
\addcontentsline{toc}{section}{Sammendrag}
Rapporten beskriver utviklingen av en lekebil som drives av et utdelt stålsvinghjul på Ø60 × 6 mm, og som skal kjøre lengst mulig etter ett opptrekk med 20 cm snor. Problemstillingen er om bilen, utviklet gjennom iterative prototyper innenfor rammene i oppgaven, kan kjøre minst \SI{5}{m} i snitt over ti kjøringer, med under \SI{15}{\percent} spredning og et sideavvik under \SI{10}{\percent} av lengden.

Tre konsepter ble vurdert i en vektet matrise: direkte drift, en trinnløs reim på kjegler og ett fast tannhjulstrinn. Tannhjulstrinnet ble valgt, med 17 : 73 ($i = 4{,}29$), Ø68-hjul og tre Ø6-aksler i 626-lagre. Den ferdige bilen veier \SI{270}{g}. Rammen er laget med generativ design i Fusion og 3D-printet i ett stykke. Lagersetet, festet av nav på aksel, snortrommelen og hjulene ble utviklet i små iterasjoner med prøvebiter. En enkel simulering av kjøringen ble brukt til å sammenligne alternativer.

Tapet i lagrene på svinghjulsakselen ble målt med utsnurring, med en Hall-sensor og en ESP32. Simuleringen viser at omtrent en firedel av energien går tapt til sluring ved nedsetting, og en like stor del i disse lagrene.

I den første kjøretesten kjørte bilen \SI[separate-uncertainty=true]{5,5 \pm 1,4}{m} i fem forsøk, med \SI{25}{\percent} spredning, og den dro til venstre i fire av dem. Det skyldes trolig at den printede rammen er vridd, slik at høyre drivhjul bærer mesteparten av lasten og tar mest av drivkraften. Distansemålet er nådd i snitt, men spredningen og sideavviket oppfyller ikke målene. De neste tiltakene er å jevne ut lasten med mellomlegg i lagerhusene.

\textbf{PLACEHOLDER: Bilde av ferdig bil}
